\guidechapter{A matrix transpose is a physical rearrangement}{Follow all nine records.}

Read the $3\times3$ grid row by row: $0,1,2,3,4,5,6,7,8$.
Now read it column by column: $0,3,6,1,4,7,2,5,8$.
The records have not changed internally. Their positions in the stream have.
The crossing paths show the exact permutation, including the records whose
positions stay fixed.

\explfigure{16-tape-order}{The geometry of one concrete data move.}{
Colors identify original rows; every numbered record is connected to its
destination. The top scan arrow represents a tape head moving through
consecutive positions. It cannot teleport to the next useful record.}

A computer program can write a new indexing formula in one line. That does
not prove that the underlying records have been physically rearranged in one
step. In the paper's machine model, a head moves at most one tape cell per
step. Large data movement must therefore be implemented and charged.

\paragraph{Why a calculation can help perform a permutation.}
For a pair of bits, three XOR updates exchange their values:
\[
 x\leftarrow x\mathbin{\mathrm{XOR}}y,\quad
 y\leftarrow y\mathbin{\mathrm{XOR}}x,\quad
 x\leftarrow x\mathbin{\mathrm{XOR}}y.
\]
For initial $(x,y)=(1,0)$ the states are $(1,0),(1,1),(0,1)$.
XOR is addition modulo two: equal bits give zero and different bits give one.
Intermediate values are combinations, yet the final result is a pure swap.
The large construction also uses intermediate linear combinations to achieve
an overall array operation. This two-bit identity only explains the possibility;
it does not establish the large construction's speed.

\paragraph{Keep two meanings of geometry separate.}
The grid here represents positions on an array. The subspaces and ranks used
later represent algebraic directions of operators on such arrays. A picture
of a short physical path is not by itself a low-rank certificate, and a
low-rank identity is not by itself a tape-time bound. The argument needs a
precise connection between them.

\takeaway{The stored order is part of the problem. A useful reordering
must be realized by actual steps, not just by renaming indices.}
